%----------------------------------------------------------------------------------------
% Tailored CV — PhD Position, Computational Immunology / IBD Genomics
% IKMB, Prof. Andre Franke, UKSH Campus Kiel — Ref. 28608
%----------------------------------------------------------------------------------------

\documentclass[10pt,a4paper,sans]{moderncv}

\usepackage[utf8]{inputenc}
\usepackage[T1]{fontenc}

\moderncvstyle{classic}
\moderncvcolor{blue}

\usepackage[scale=0.78]{geometry}
\usepackage{ragged2e}

%----------------------------------------------------------------------------------------

\firstname{Kundai Farai}
\familyname{Sachikonye}

\address{Auf der Lichtung 19}{82178 Puchheim, Germany}
\mobile{(+49) 1626386344}
\email{kundai.sachikonye@bitspark.com}
\homepage{github.com/fullscreen-triangle}{github.com/fullscreen-triangle}

%----------------------------------------------------------------------------------------

\begin{document}

\makecvtitle

%----------------------------------------------------------------------------------------
%	EDUCATION
%----------------------------------------------------------------------------------------

\section{Education}

\cventry{2019--2021}{PhD candidate, Computational Lipidomics}{Technische Universität München / Bayerisches Forschungsinstitut}{}{}{
  ML for large-scale omics data in clinical cohorts, multi-omics statistical
  modelling, distributed inference across federated clinical sites.
  Departed to pursue independent research.
}
\cventry{2018--2019}{MSc Computational Biology}{Universität Tübingen}{}{}{
  Probabilistic graphical models, statistical bioinformatics, Bayesian methods,
  computational systems biology, ML for biological sequences.
}
\cventry{2014--2017}{MSc Pharmaceutical Biotechnology}{HAW Hamburg}{}{}{}
\cventry{2011--2014}{BSc Biochemistry and Cell Biology}{Jacobs University Bremen}{}{}{}

%----------------------------------------------------------------------------------------
%	RELEVANT MANUSCRIPTS
%----------------------------------------------------------------------------------------

\section{Relevant Manuscripts}

\cvitem{2025}{
  \textbf{On the Thermodynamic Consequences of Categorical Completion on Biological Systems.}
  Spectral complementarity $\rho$ for MHC-I/II binding: AUC 0.81 vs.\ NetMHCpan 0.68
  ($r=0.73$, $p<10^{-12}$, $n=2{,}847$ IEDB peptides).
  Neoantigen identification from dual $\rho$ threshold. \textit{Manuscript.}
}
\cvitem{2025}{
  \textbf{On Disease Epistemology in Blind Receiver.}
  Hierarchical Bayesian disease state $\mathbf{D}\in[0,1]^8$; AUC$=1.00$;
  25/25 validation experiments; sparse $\ell_1$ LP for personalised therapeutic
  target selection. \textit{Manuscript. AIMe Registry / TUM Weihenstephan.}
}

%----------------------------------------------------------------------------------------
%	COMPUTATIONAL IMMUNOLOGY AND GENOMICS TOOLS
%----------------------------------------------------------------------------------------

\section{Computational Immunology and Genomics Tools}

\cvitem{syndrome}{
  \textbf{Immunopeptidology Toolkit for IBD Immune Analysis.}
  3-channel DFT (hydropathy, volume, charge) $\to$ 36D cosine similarity $\rho$;
  $\mathcal{O}(cL\log L)$; no training data.
  \textbf{MHC binding predictor:} $\rho\geq0.72$; AUC 0.81 vs.\ NetMHCpan 0.68
  on 2,847 IEDB peptides. Patient HLA scanning: predict which microbial peptides
  from metagenomics are presented and drive mucosal T cell activation.
  \textbf{Cross-reactive epitope screen:} dual threshold
  $\Delta\rho_\mathrm{self}>0.15$ + $\rho_\mathrm{MHC}>0.72$ for molecular
  mimicry candidates (microbial epitopes cross-reactive to intestinal self-antigens).
  \textbf{Escape predictor:} $E(v)=-\max\rho(v,\mathrm{self})+w\cdot\rho(v,\mathrm{target})$;
  maps TCR variants maintaining target recognition while minimising autoreactive risk.
  \textbf{Disease state layer:} $\mathbf{D}\in[0,1]^8$ including $D_\mathrm{Imm}$
  (immune activation), $D_\mathrm{Bar}$ (mucosal barrier), $D_\mathrm{ATP}$,
  $D_\mathrm{Circ}$; sparse $\ell_1$ LP for minimum-intervention target selection.
  Live: \href{https://syndrome-seven.vercel.app/tools}{syndrome-seven.vercel.app/tools}.
  \href{https://github.com/fullscreen-triangle/syndrome}{github.com/fullscreen-triangle/syndrome}
}

\cvitem{gospel}{
  \textbf{Large-Scale Genomics Pipeline for IBD Variant Analysis.}
  $10^6$ variants/second VCF processing; germline and somatic variant calling;
  RNA-seq quantification and differential expression; Bayesian network inference
  for multi-omics integration (host genotype $\to$ metagenomics $\to$ transcriptomics
  as conditional graph); GMMs for population structure.
  Validated: 1000~Genomes Phase~3, gnomAD~v3.1.2, UK~Biobank.
  ClinVar, KEGG, Reactome API integration; Ray distributed; Docker deployable.
  \href{https://github.com/fullscreen-triangle/gospel}{github.com/fullscreen-triangle/gospel}
}

\cvitem{four-sided-triangle}{
  \textbf{HPC Inference Pipeline with Automated Evaluation.}
  Rust core (15--50$\times$ speedup via PyO3 bindings); Ray-distributed parallel
  execution; FastAPI orchestration; Docker/Kubernetes.
  8-stage metacognitive pipeline with automated benchmarking at each stage.
  \href{https://github.com/fullscreen-triangle/four-sided-triangle}{github.com/fullscreen-triangle/four-sided-triangle}
}

\cvitem{lavoisier}{
  \textbf{Multi-Omics Data Pipeline.}
  Dual pipeline: numerical MS1/MS2 processing + CNN spectral classification (PyTorch);
  98.9\,\% NIST accuracy; Ray distributed; $>$100\,GB mzML; memory-mapped I/O;
  HDF5 storage. Applicable to metagenomics data preprocessing pipelines at cohort scale.
  \href{https://github.com/fullscreen-triangle/lavoisier}{github.com/fullscreen-triangle/lavoisier}
}

%----------------------------------------------------------------------------------------
%	EXPERIENCE
%----------------------------------------------------------------------------------------

\section{Experience}

\cventry{2021--present}{Independent AI \& Computational Biology Developer}{Bitspark GmbH}{}{}{
  Computational immunopeptidology, large-scale genomics, Bayesian multi-omics
  integration, and HPC pipeline development. All systems publicly deployed and validated.
}

\cventry{2019--2021}{Computational Lipidomics}{TUM / Bayerisches Forschungsinstitut}{Munich}{}{
  ML for clinical omics cohorts: CNN spectral classification, distributed pipeline
  maintenance, multi-omics integration across federated clinical sites. Python, Java.
}

\cventry{2017--2018}{Glycomics and Proteomics}{Universitätsklinikum Hamburg-Eppendorf}{}{}{
  Automated LC-MS/MS and MALDI-TOF pipeline; clinical sample processing;
  data annotation and quality control for clinical cohort studies.
}

%----------------------------------------------------------------------------------------
%	TECHNICAL SKILLS
%----------------------------------------------------------------------------------------

\section{Technical Skills}

\cvitem{Computational Immunology}{
  MHC binding prediction (spectral DFT, validated vs.\ NetMHCpan);
  T cell epitope identification; cross-reactive epitope screening;
  TCR/BCR sequence analysis; V(D)J combinatorial statistics;
  immune repertoire distance metrics (spectral $\rho$ clustering);
  neoantigen identification; immune escape scoring
}

\cvitem{Genomics / Multi-Omics}{
  GWAS-scale variant analysis (gnomAD, UK~Biobank, 1000~Genomes validated);
  RNA-seq quantification and differential expression; metagenomics pipeline
  integration; Bayesian network multi-omics integration (host genotype,
  metagenomics, transcriptomics joint inference); population structure (GMMs)
}

\cvitem{Bayesian Methods / Statistics}{
  Variational Bayes (ELBO optimisation); hierarchical / multi-level models;
  probabilistic graphical models; Bayesian network inference;
  sparse $\ell_1$ regularisation; Gaussian processes; MCMC (NUTS/HMC);
  continuous shrinkage priors; automatic differentiation (PyTorch \texttt{autograd})
}

\cvitem{Single-Cell / Sequencing}{
  Scanpy / Seurat pipelines (UMAP, Leiden, diffusion pseudotime);
  RNA velocity trajectory inference; probabilistic dimensionality reduction;
  HDF5 / AnnData; FASTQ, BAM, VCF data structures; samtools; Nextflow
}

\cvitem{Languages / Infrastructure}{
  Python (primary), R, Rust, Java, TypeScript, Bash;
  Ray distributed; Docker, Kubernetes; memory-mapped I/O;
  gnomAD, ClinVar, KEGG, Reactome APIs; Git, CI/CD
}

%----------------------------------------------------------------------------------------
%	LANGUAGES
%----------------------------------------------------------------------------------------

\section{Languages}

\cvitemwithcomment{English}{Native / Fluent}{}
\cvitemwithcomment{German}{Conversational}{Living in Germany since 2011}

%----------------------------------------------------------------------------------------

\renewcommand{\listitemsymbol}{-~}

\end{document}
